\documentclass[11pt,a4paper]{article}
\usepackage[utf8]{inputenc}
\usepackage[T1]{fontenc}
\usepackage{lmodern}
\usepackage[french]{babel}
\usepackage{amsmath,amssymb,mathtools}
\usepackage{icomma}
\usepackage[version=4]{mhchem}
\usepackage{siunitx}
\sisetup{output-decimal-marker={,},per-mode=symbol}
\usepackage{xcolor}
\usepackage{colortbl}
\usepackage{tikz}
\usepackage[most]{tcolorbox}
\usepackage{enumitem}
\usepackage{array}
\usepackage{booktabs}
\usepackage{fancyhdr}
\usepackage[hidelinks]{hyperref}
\usetikzlibrary{arrows.meta,positioning,calc}
\usepackage[a4paper,margin=2.1cm,headheight=26pt,headsep=10pt]{geometry}

\definecolor{primary}{RGB}{146,95,161}
\definecolor{primarydark}{RGB}{92,55,108}
\definecolor{acidecol}{RGB}{205,60,45}
\definecolor{basecol}{RGB}{40,110,200}
\definecolor{excol}{RGB}{0,135,90}
\definecolor{attncol}{RGB}{200,40,55}

\tcbset{boxbase/.style={enhanced,breakable,boxrule=0.9pt,arc=2.5pt,left=3mm,right=3mm,top=2.3mm,bottom=2.3mm,fonttitle=\bfseries,coltitle=white,toptitle=0.6mm,bottomtitle=0.6mm}}
\newtcolorbox[auto counter]{corrbox}[1][]{boxbase,colback=excol!4,colframe=excol,title={\textsc{Corrigé}~\thetcbcounter\ifx\relax#1\relax\relax\else\textmd{ --- \itshape #1}\fi}}
\newtcolorbox{astuce}{boxbase,colback=attncol!6,colframe=attncol,sharp corners}

\pagestyle{fancy}
\fancyhf{}
\fancyhead[L]{\small\textcolor{primarydark}{\textbf{Fiche 10} — Thème 2 : pH, pOH, $K_e$}}
\fancyhead[R]{\small\textcolor{primarydark}{\textsc{Corrigé — Facile}}}
\fancyfoot[C]{\small\thepage}
\renewcommand{\headrulewidth}{0.8pt}
\renewcommand{\headrule}{\hbox to\headwidth{\color{primary}\leaders\hrule height \headrulewidth\hfill}}

\setlength{\parindent}{0pt}
\setlength{\parskip}{3pt}

\begin{document}

\begin{center}
{\LARGE\bfseries\color{primarydark} Thème 2 --- Corrigé Facile}
\end{center}

\begin{corrbox}[du pH vers la concentration]
On applique $[\ce{H3O+}] = 10^{-\mathrm{pH}}$ :
\begin{enumerate}[label=\alph*),leftmargin=1.6em,itemsep=1pt]
  \item $[\ce{H3O+}] = 10^{-1} = 0{,}1\ \si{mol.L^{-1}}$
  \item $[\ce{H3O+}] = 10^{-4}\ \si{mol.L^{-1}}$
  \item $[\ce{H3O+}] = 10^{-2}\ \si{mol.L^{-1}}$
  \item $[\ce{H3O+}] = 10^{-7}\ \si{mol.L^{-1}}$ (eau neutre)
\end{enumerate}
\end{corrbox}

\begin{corrbox}[de la concentration vers le pH]
On applique $\mathrm{pH} = -\log[\ce{H3O+}]$ :
\begin{enumerate}[label=\alph*),leftmargin=1.6em,itemsep=1pt]
  \item $\mathrm{pH} = -\log(10^{-3}) = 3$
  \item $\mathrm{pH} = -\log(10^{-9}) = 9$
  \item $\mathrm{pH} = -\log(5\times10^{-4}) = 4 - \log 5 \approx 4 - 0{,}70 = \mathbf{3{,}30}$
  \item $\mathrm{pH} = -\log(2\times10^{-6}) = 6 - \log 2 \approx 6 - 0{,}30 = \mathbf{5{,}70}$
\end{enumerate}
\end{corrbox}

\begin{corrbox}[acide, neutre ou basique ?]
Frontière à $\mathrm{pH}=7$ (soit $[\ce{H3O+}]=[\ce{OH-}]=10^{-7}$) :
\begin{enumerate}[label=\alph*),leftmargin=1.6em,itemsep=1pt]
  \item $\mathrm{pH}=3{,}5<7$ : \textcolor{acidecol}{\textbf{acide}}.
  \item $\mathrm{pH}=7$ : \textbf{neutre}.
  \item $\mathrm{pH}=10{,}2>7$ : \textcolor{basecol}{\textbf{basique}}.
  \item $[\ce{H3O+}]=10^{-5}>10^{-7}$ : \textcolor{acidecol}{\textbf{acide}} (\,$\mathrm{pH}=5$\,).
  \item $[\ce{OH-}]=10^{-3}$, donc $[\ce{H3O+}]=10^{-14}/10^{-3}=10^{-11}<10^{-7}$ :
        \textcolor{basecol}{\textbf{basique}} (\,$\mathrm{pH}=11$\,).
\end{enumerate}
\end{corrbox}

\begin{corrbox}[passer par $\mathrm{pH}+\mathrm{pOH}=14$]
$\mathrm{pOH}=14-\mathrm{pH}$ puis $[\ce{OH-}]=10^{-\mathrm{pOH}}$ :
\begin{enumerate}[label=\alph*),leftmargin=1.6em,itemsep=1pt]
  \item $\mathrm{pOH}=11$, $[\ce{OH-}]=10^{-11}\ \si{mol.L^{-1}}$.
  \item $\mathrm{pOH}=3$, $[\ce{OH-}]=10^{-3}\ \si{mol.L^{-1}}$.
  \item $\mathrm{pOH}=5{,}5$, $[\ce{OH-}]=10^{-5{,}5}\approx 3{,}2\times10^{-6}\ \si{mol.L^{-1}}$.
\end{enumerate}
\end{corrbox}

\begin{corrbox}[le produit ionique en action]
\begin{enumerate}[label=\alph*),leftmargin=1.6em,itemsep=1pt]
  \item $[\ce{OH-}] = \dfrac{10^{-14}}{10^{-2}} = 10^{-12}\ \si{mol.L^{-1}}$.
  \item $[\ce{H3O+}] = \dfrac{10^{-14}}{10^{-4}} = 10^{-10}\ \si{mol.L^{-1}}$.
  \item $[\ce{H3O+}]^2 = K_e \Rightarrow [\ce{H3O+}] = \sqrt{10^{-14}} = 10^{-7}\ \si{mol.L^{-1}}$
        (et $[\ce{OH-}]=10^{-7}$).
\end{enumerate}
\end{corrbox}

\begin{astuce}
\textbf{Vérification de cohérence.} Après tout calcul, contrôlez que
$[\ce{H3O+}]\times[\ce{OH-}] = 10^{-14}$ et que $\mathrm{pH}+\mathrm{pOH}=14$. Si ce n'est pas le cas, il y a
une erreur de signe ou d'exposant.
\end{astuce}

\end{document}
